\songtitle{Don't Mistake the Bridge for Home}

\tag*{2026}\tag{folk-rock}\tag{latin-jazz}\tag{original-source}\tag{thomas-entourage}\tag{thomas-garcia}\tag{zahra-moradi}\label{sec:dont-mistake-the-bridge-for-home}
\begin{notes}
Lyrics by Suno based on \emph{Lines Not Crossed} by Carlos Thompson, from the \emph{Thomas' Entourage} narrative universe, a duet of Thomas and Zahra during her visit to Punta Espada in December 2026 -- the drive from the villa between desert and sea, a session where the pütchipü'üi still carries the words between the parties, and what progress ought to mean. From the same story as \emph{Stopped Before the Edge} (page \pageref{sec:stopped-before-the-edge}), \emph{Where the Water Meets Me} (page \pageref{sec:where-the-water-meets-me}) and \emph{Keep Your Hands Kind} (page \pageref{sec:keep-your-hands-kind}). Retrieved from Gabriel's Suno page (V6-MINI, 3 November 2026).
\end{notes}
\begin{style}
Folk-rock fused with Latin-jazz instrumentation: steady roots-rock pulse, nylon-string guitar, warm electric bass, brushed kit, hand percussion, piano montuno accents, muted trumpet and agile Latin-jazz piano fills, with fiddle threading the folk melody. Intimate narrative verses trade between two distinct duet voices, then open into close harmonies on a memorable chorus; conversational, observant, quietly hopeful. Organic live-room feel, earthy tape warmth, clear vocals, dynamic arrangement that builds gently through the afternoon-session verse.
\end{style}
\begin{lyrics}
\songpart{Verse 1: Thomas}
The road ran quiet from the villa gate,\\
You watched my hands turn through the bends.\\
We used to rent a car in Baltic towns,\\
And let the afternoon decide where it ends.

\songpart{Verse 2: Zahra}
You could have had a driver waiting there,\\
But chose the wheel and this small stretch of time.\\
Desert on one side, sea on the other,\\
No need to fill the silence with a line.

\songpart{Pre-Chorus: Duet}
A little room between the day and all its demands,\\
A little time to see the land.

\songpart{Chorus: Duet}
Let the tools make the long road shorter,\\
Let the old words keep their place.\\
Build a bridge that carries every voice,\\
But don’t mistake the bridge for home.\\
When the tables turn and the hours run on,\\
We can change the way the work gets done—\\
Still, the people have to lead the way.

\songpart{Verse 3: Zahra}
The session moved through forms and careful phrases,\\
Then one sharp question set the room alight.\\
I wondered what the interpreter softened,\\
What bright-edged words went missing in the flight.

\songpart{Verse 4: Thomas}
Screens linked the scattered hands together,\\
Fewer steps between a thought and its reply.\\
Still the pütchipü'üi carried words between us,\\
A living trust no network can supply.

\songpart{Pre-Chorus: Duet}
At lunch we said what progress ought to mean:\\
Make room for lives, don’t take their place between.

\songpart{Chorus: Duet}
Let the tools make the long road shorter,\\
Let the old words keep their place.\\
Build a bridge that carries every voice,\\
But don’t mistake the bridge for home.\\
When the tables turn and the hours run on,\\
We can change the way the work gets done—\\
Still, the people have to lead the way.

\songpart{Bridge: Duet}
The morning’s questions, the afternoon’s fire,\\
A steady hand and a rising reply.\\
The smoothest day can still hold a storm;\\
Tradition stays alive when it can transform.

\songpart{Final Chorus: Duet}
Let the tools make the long road shorter,\\
Let the old words keep their place.\\
Build a bridge that carries every voice,\\
But don’t mistake the bridge for home.\\
When the tables turn and the hours run on,\\
We can change the way the work gets done—\\
Still, the people have to lead the way.

\songpart{Outro: Duet}
Sea beside us, desert wide,\\
A quiet road, two voices side by side.
\end{lyrics}

